\documentclass[11pt,a4paper]{article}
\usepackage[utf8]{vietnam}
\usepackage{amsmath,amssymb}
\usepackage{tabularx}
\usepackage{tikz}
\usepackage{xcolor}
\usepackage{fancyhdr}
\usepackage[top=1.8cm, bottom=2.0cm, left=1.8cm, right=1.8cm]{geometry}

\setlength{\parskip}{3pt}
\setlength{\parindent}{0pt}

\pagestyle{fancy}
\fancyhf{}
\renewcommand{\headrulewidth}{0.4pt}
\lhead{\footnotesize \textbf{CHÙA HƯƠNG SƠN -- HƯỚNG DẪN CHẤM THI}}
\rhead{\footnotesize \textbf{ĐÁP ÁN CHI TIẾT ĐỀ THI THỬ SỐ 05}}
\cfoot{\thepage}

\begin{document}

\begin{center}
{\large\bfseries ĐÁP ÁN \& HƯỚNG DẪN CHẤM ĐỀ THI THỬ CUỐI HỌC KỲ I (ĐỀ SỐ 05)}\\[1.5mm]
{\bfseries ĐƠN VỊ: CHÙA HƯƠNG SƠN \quad|\quad MÔN: TOÁN -- LỚP 9 \quad|\quad NĂM HỌC 2026 -- 2027}\\[1mm]
\textit{(Chuẩn cấu trúc định dạng mới của Sở Giáo dục và Đào tạo)}
\end{center}

\vspace{1mm}
\hrule
\vspace{3mm}

\noindent
\textbf{\large PHẦN I. TRẮC NGHIỆM KHÁCH QUAN (7,0 ĐIỂM)}

\vspace{2mm}
\noindent
\textbf{1. Bảng đáp án Phần 1: Trắc nghiệm nhiều lựa chọn (3,0 điểm)} \textit{(Mỗi câu đúng được 0,25 điểm)}
\begin{center}
\begin{tabular}{|c|c|c|c|c|c|c|c|c|c|c|c|c|}
\hline
\textbf{Câu} & \textbf{1} & \textbf{2} & \textbf{3} & \textbf{4} & \textbf{5} & \textbf{6} & \textbf{7} & \textbf{8} & \textbf{9} & \textbf{10} & \textbf{11} & \textbf{12} \\
\hline
\textbf{Đáp án} & \textbf{A} & \textbf{A} & \textbf{A} & \textbf{A} & \textbf{A} & \textbf{A} & \textbf{A} & \textbf{A} & \textbf{A} & \textbf{A} & \textbf{A} & \textbf{A} \\
\hline
\end{tabular}
\end{center}

\vspace{2mm}
\noindent
\textbf{2. Bảng đáp án Phần 2: Trắc nghiệm Đúng -- Sai (2,0 điểm)}
\begin{center}
\begin{tabular}{|c|c|c|c|c|l|}
\hline
\textbf{Câu} & \textbf{Ý a} & \textbf{Ý b} & \textbf{Ý c} & \textbf{Ý d} & \textbf{Ghi chú đáp án chuẩn} \\
\hline
\textbf{Câu 13} & \textbf{Đ} & \textbf{Đ} & \textbf{S} & \textbf{Đ} & Ý c) sai vì lượng nước tháng này bằng $74\text{ m}^3$ (đề ghi $80$). \\
\hline
\textbf{Câu 14} & \textbf{Đ} & \textbf{Đ} & \textbf{Đ} & \textbf{S} & Ý d) sai vì $OA \perp EF$ (chứ không song song). \\
\hline
\end{tabular}
\end{center}

\vspace{2mm}
\noindent
\textbf{3. Bảng đáp án Phần 3: Trắc nghiệm trả lời ngắn (2,0 điểm)} \textit{(Mỗi câu đúng được 0,5 điểm)}
\begin{center}
\begin{tabular}{|c|l|c|c|}
\hline
\textbf{Câu} & \textbf{Nội dung câu hỏi} & \textbf{Kết quả bài làm} & \textbf{Điểm} \\
\hline
\textbf{Câu 15} & Nghiệm bất phương trình $3(2x - 1) + 4 > 4(x + 2) - 1$ & \textbf{$x > 3$} & 0,5đ \\
\hline
\textbf{Câu 16} & Giá trị biểu thức $H = \sqrt{(\sqrt{11}-4)^2} + \sqrt{11}$ & \textbf{$4$} & 0,5đ \\
\hline
\textbf{Câu 17} & Chiều cao cây trước khi gãy ($6 \cdot \tan 35^\circ + \dfrac{6}{\cos 35^\circ}$) & \textbf{$11{,}52\text{ m}$} & 0,5đ \\
\hline
\textbf{Câu 18} & Diện tích hình quạt tròn ($\dfrac{3{,}14 \times 12^2 \times 60}{360}$) & \textbf{$75{,}4\text{ cm}^2$} & 0,5đ \\
\hline
\end{tabular}
\end{center}

\vspace{3mm}
\noindent
\textbf{\large PHẦN II. TỰ LUẬN (3,0 ĐIỂM)}

\vspace{2mm}
\textbf{Câu 19 (1,0 điểm).}
\begin{itemize}
  \item[a)] \textit{(0,5 điểm)} Cộng hai phương trình: $4x = 8 \Leftrightarrow x = 2$. Thay vào phương trình thứ hai: $2 - 2y = -4 \Leftrightarrow 2y = 6 \Leftrightarrow y = 3$. Hệ có nghiệm duy nhất: $(x; y) = (2; 3)$. \hfill \textbf{(0,5đ)}
  \item[b)] \textit{(0,5 điểm)} Quy đồng: $Q = \dfrac{\sqrt{x}-2+\sqrt{x}+2}{(\sqrt{x}+2)(\sqrt{x}-2)} \cdot \dfrac{x-4}{2\sqrt{x}} = \dfrac{2\sqrt{x}}{x-4} \cdot \dfrac{x-4}{2\sqrt{x}} = 1$. \hfill \textbf{(0,5đ)}
\end{itemize}

\vspace{1mm}
\textbf{Câu 20 (1,5 điểm).}
\begin{itemize}
  \item[a)] \textit{(0,75 điểm)} Vì $AD \perp BC \Rightarrow \widehat{ADH} = 90^\circ$; $BE \perp AC \Rightarrow \widehat{AEH} = 90^\circ$. Gọi $I$ là trung điểm $AH \Rightarrow IA = IE = IH = ID = \dfrac{AH}{2} \Rightarrow A, E, H, D$ cùng nằm trên đường tròn đường kính $AH$. \hfill \textbf{(0,75đ)}
  \item[b)] \textit{(0,75 điểm)} Vì $AK$ là đường kính của $(O) \Rightarrow \widehat{ABK} = 90^\circ \Rightarrow BK \perp AB$. Lại có $CH \perp AB \Rightarrow CH \parallel BK$. Tương tự: $\widehat{ACK} = 90^\circ \Rightarrow CK \perp AC$. Lại có $BH \perp AC \Rightarrow BH \parallel CK$. Tứ giác $BHCK$ có các cặp cạnh đối song song nên là hình bình hành (đpcm). \hfill \textbf{(0,75đ)}
\end{itemize}

\vspace{1mm}
\textbf{Câu 21 (0,5 điểm).}
\begin{itemize}
  \item[a)] \textit{(0,25 điểm)} Thể tích $V = \pi R^2 h = 16\pi \Rightarrow h = \dfrac{16}{R^2}$. Diện tích toàn phần:
  $S_{tp} = 2\pi R^2 + 2\pi R h = 2\pi R^2 + \dfrac{32\pi}{R} = 2\pi R^2 + \dfrac{16\pi}{R} + \dfrac{16\pi}{R} \ge 3\sqrt[3]{2\pi R^2 \cdot \dfrac{16\pi}{R} \cdot \dfrac{16\pi}{R}} = 3\sqrt[3]{512\pi^3} = 24\pi\text{ m}^2$. Dấu ``='' xảy ra khi $2\pi R^2 = \dfrac{16\pi}{R} \Leftrightarrow R^3 = 8 \Leftrightarrow R = 2\text{ m} \Rightarrow h = 4\text{ m}$. \hfill \textbf{(0,25đ)}
  \item[b)] \textit{(0,25 điểm)} Kỹ thuật Cauchy ngược dấu: $\dfrac{a}{b^2+1} = a - \dfrac{ab^2}{b^2+1} \ge a - \dfrac{ab^2}{2b} = a - \dfrac{ab}{2}$. Tương tự: $\dfrac{b}{c^2+1} \ge b - \dfrac{bc}{2}$ và $\dfrac{c}{a^2+1} \ge c - \dfrac{ca}{2}$. Cộng lại: $P \ge (a+b+c) - \dfrac{ab+bc+ca}{2} \ge \sqrt{3(ab+bc+ca)} - \dfrac{3}{2} = 3 - \dfrac{3}{2} = \dfrac{3}{2}$. Dấu ``='' xảy ra khi $a = b = c = 1$. Vậy $\min P = \dfrac{3}{2}$. \hfill \textbf{(0,25đ)}
\end{itemize}

\end{document}
